\scanpage{97}
% archive foliation: 96 No author page-number visible.
comme
\[
(62)\quad\left\{\begin{aligned}
T_{\Sigma_J,t}&\simeq O_t(2)(\underline\omega)\\
\text{avec }O_t(2)&\overset{\mathrm{def}}\simeq O_t(1)^{\otimes2},\qquad
O_t(1)\simeq V_J/F_t\simeq\check F_t(\underline\omega)
\end{aligned}\right.
\]
de sorte que $\nu_E$ s'interprète comme un isom
\[
(63)\qquad\nu:\underline t_{\unsure{E}}^{\otimes2}\simeq
\deleted{\unsure{\check F}}F_t^{\otimes(-2)}(\underline\omega)
\]
% margin position: right
\margin{ou aussi comme $\check\nu:\underline t^{\otimes-2}\simeq F_t^{\otimes2}(\underline\omega)$}
(NB $\underline\omega=\underline\omega_{\underline J}$, ens. des deux ordres circulaires
sur $\underline J$). Les données de base sont donc
pour nous
\[
\left[J\ (\mapsto\underline\omega=\underline\omega_J\unsure{;}\,,\ \unsure{t}\text{ et }V_J=V_J(k)),\quad
F_t\hookrightarrow V_J,\quad\underline t,\quad
\nu:\underline t^{\otimes2}\simeq F_t^{\otimes(-2)}(\underline\omega)\right].
\]
Ceci posé, la donnée d'un épinglage de
Legendre s'explicite comme un isom
\[
(64)\qquad\kappa:\underline t\xrightarrow{\sim}\underline O_t(1)^{\unsure{\sharp}}
\]
i.e. comme une famille d'isomorphismes
\[
(65)\quad\left\{\begin{aligned}
&(\kappa_i)_{i\in\mu_4^*},\qquad\kappa_i:\underline t\xrightarrow{\sim}\underline O_t(1)\simeq\check F_t(\underline\omega)\\
&\kappa_{i'}=i\kappa_i\qquad\text{où }i'=1/i\text{ est l'autre élément de }\mu_4^*
\end{aligned}\right.
\]
% margin position: left
\margin{$\mu_4^*=\unsure{\boldsymbol{\mu}|\mu}_4^*(k)$, éléments d'ordre $4$ de $k^*$.}
ou encore, d'une famille
\[
(66)\quad\left\{\begin{aligned}
&(\check\kappa_i)_{i\in\mu_4^*},\qquad\check\kappa_i:\check{\underline t}\simeq F_t(\underline\omega)\\
&\check\kappa_{i'}=\frac1i\check\kappa_i,
\end{aligned}\right.
\]
satisfaisant à une certaine condition:
\[
(67)\quad\left\{\begin{aligned}
&\kappa_i^{\otimes2}:\underline t^{\otimes2}\xrightarrow{\sim}\unsure{F_t|\check F_t}^{\otimes(-2)}\quad\text{est de la forme}\\
&\kappa_i^{\otimes2}(\delta^{\otimes2})=\nu(\delta^{\otimes2})\wedge\omega_i.
\end{aligned}\right.
\]
où $\omega_i\in\underline\omega$ est un élément convenable de $\underline\omega$,
bien déterminé bien sûr par $\kappa,i$, en fait par $\kappa_i$.
